\usepackage[T1]{fontenc}
\usepackage[utf8]{inputenc}
\usepackage{lmodern}
\usepackage{amsmath,amssymb,mathtools}
\usepackage{tikz}
\usepackage[margin=1.15in]{geometry}
\usepackage[final]{microtype}
\usepackage[colorlinks=true,linkcolor=blue,citecolor=blue,urlcolor=blue]{hyperref}
\usepackage{enumitem}

\newtheorem{theorem}{Theorem}[section]
\newtheorem{lemma}[theorem]{Lemma}
\newtheorem{proposition}[theorem]{Proposition}
\newtheorem{corollary}[theorem]{Corollary}
\newtheorem*{claim}{Claim}
\theoremstyle{definition}
\newtheorem{definition}[theorem]{Definition}
\theoremstyle{remark}
\newtheorem{remark}[theorem]{Remark}
\newtheorem{question}[theorem]{Question}

\numberwithin{equation}{section}

\newcommand{\rr}{\mathfrak{rr}}
\newcommand{\rrp}{\mathfrak{rr}'}
\newcommand{\rrf}{\mathfrak{rr}_{f}}
\newcommand{\rri}{\mathfrak{rr}_{i}}
\newcommand{\rro}{\mathfrak{rr}_{o}}
\newcommand{\rrfi}{\mathfrak{rr}_{fi}}
\newcommand{\rrfo}{\mathfrak{rr}_{fo}}
\newcommand{\rrio}{\mathfrak{rr}_{io}}
\newcommand{\sr}{\mathfrak{sr}}
\newcommand{\sz}{\mathfrak{sz}}
\newcommand{\bb}{\mathfrak b}
\newcommand{\dd}{\mathfrak d}
\newcommand{\reap}{\mathfrak r}
\newcommand{\linpred}{\mathfrak v_{\ell}}
\newcommand{\cc}{\mathfrak c}
\newcommand{\ddensity}{\mathfrak{dd}}
\newcommand{\cM}{\mathcal M}
\newcommand{\cN}{\mathcal N}
\newcommand{\covN}{\operatorname{cov}(\cN)}
\newcommand{\cofM}{\operatorname{cof}(\cM)}
\newcommand{\cofN}{\operatorname{cof}(\cN)}
\newcommand{\nonM}{\operatorname{non}(\cM)}
\newcommand{\Sym}{\operatorname{Sym}(\omega)}
\newcommand{\CC}{\mathsf{CC}}
\newcommand{\E}{\mathbb E}
\newcommand{\PP}{\mathbb P}
\newcommand{\QBS}{\mathbb Q_{\mathrm{BS}}}
\newcommand{\Qfg}{\mathbb Q_{f,g}}
\newcommand{\Pow}{\mathcal P}
\newcommand{\Dp}{\operatorname{Dp}}
\newcommand{\girth}{\operatorname{girth}}
\newcommand{\restr}{\mathbin{\upharpoonright}}
\newcommand{\abs}[1]{\left| #1\right|}
\newcommand{\dom}{\operatorname{dom}}
\newcommand{\ran}{\operatorname{ran}}

\title{Determining the Rearrangement Number}
\author{Vinicius de Oliveira Rodrigues}
\address{University of São Paulo, Rua do Matão, 1010, São Paulo, SP, Brazil}
\email{vinior@ime.usp.br}
\date{}

\hypersetup{
  pdftitle={Determining the Rearrangement Number},
  pdfauthor={Vinicius de Oliveira Rodrigues},
  pdfsubject={The rearrangement number and the uniformity of the meager ideal},
  pdfkeywords={rearrangement number, cardinal characteristics, slaloms, conditionally convergent series}
}

\subjclass[2020]{Primary 03E17; Secondary 40A05}
\keywords{rearrangement number, cardinal characteristics, slaloms, conditionally convergent series}
